\documentclass[11pt]{article}
\usepackage[margin=1in]{geometry}
\usepackage{amsmath,amssymb,amsthm,hyperref}
\newtheorem{theorem}{Theorem}
\newtheorem{proposition}{Proposition}
\newcommand{\Gq}{\Gamma_q}
\title{A negative-parameter counterexample to the\\
unrestricted q-gamma extension\\
\large MathDB 335285}
\author{Principia Math in Progress}
\date{12 September 2026}
\begin{document}
\maketitle

\paragraph{Task 1: resolve the stated extension.}
Conjecture 1 of Kalmykov--Karp \cite[p.~344]{KK} proposes extending
all conclusions of their Theorem 3 to every $\alpha>0$, retaining
the other hypotheses. For a nontrivial nonnegative log-concave sequence
$(g_n)$ with no internal zeros, that theorem concerns
\[
 g(\mu;x)=\sum_{n\ge0}\frac{g_nx^n}{\Gq(\mu+n)},\qquad
 \Delta_g=g(\mu+\alpha;x)g(\mu+\beta;x)
             -g(\mu;x)g(\mu+\alpha+\beta;x).
\]
The retained parameter condition is
$\beta>0$ and $\mu\ge\max(-\beta,-1)$
\cite[p.~342]{KK}.
The proposed positivity is false on this full domain.
The counterexample below does not decide the separately restricted
problem with $\mu\ge0$.

\begin{theorem}
Set
\[
 q=\frac{256}{625},\qquad
 \mu=-\frac34,\qquad \alpha=\frac14,\qquad \beta=2,
 \qquad g_0=1,\quad g_n=0\quad(n\ge1).
\]
This sequence satisfies the hypotheses of Theorem 3, its series has
radius of convergence infinity, and for every $x>0$,
\[
 \boxed{\displaystyle
 \Delta_g=-\frac{95}{576\,\Gq(1/4)\Gq(1/2)}<0.}
\]
In particular, both the value-positivity and coefficient-positivity
conclusions of the proposed unrestricted extension fail.
\end{theorem}

\begin{proof}
For $0<q<1$, write $[z]_q=(1-q^z)/(1-q)$.
The defining product
\[
 \Gq(z)=(1-q)^{1-z}\frac{(q;q)_\infty}{(q^z;q)_\infty},
 \qquad (a;q)_\infty=\prod_{j\ge0}(1-aq^j),
\]
gives $\Gq(z+1)=[z]_q\Gq(z)$ wherever both sides are finite.
For positive real $z$ its product is positive and finite: every factor
is positive, and the geometric tails have finite sums.
Put $A=1/\Gq(1/4)>0$ and $B=1/\Gq(1/2)>0$.
The four arguments in $\Delta_g$ are
$-3/4,-1/2,5/4,3/2$, so none is a pole.
The recurrence gives
\[
 g(-3/4;x)=[-3/4]_q A,\quad
 g(-1/2;x)=[-1/2]_q B,\quad
 g(5/4;x)=\frac{A}{[1/4]_q},\quad
 g(3/2;x)=\frac{B}{[1/2]_q}.
\]
Consequently,
\[
 \frac{\Delta_g}{AB}
 =\frac{[-1/2]_q}{[1/4]_q}
  -\frac{[-3/4]_q}{[1/2]_q}
 =\frac{1-25/16}{1-4/5}
  -\frac{1-125/64}{1-16/25}
 =-\frac{45}{16}+\frac{1525}{576}
 =-\frac{95}{576}.
\]
Here $q^{1/4}=4/5$ was used. The sequence has support $\{0\}$,
is nontrivial, and satisfies $g_n^2\ge g_{n-1}g_{n+1}$ for every
$n\ge1$; its zeros are all terminal. Moreover,
$\mu=-3/4>-1=\max(-2,-1)$ and $\alpha>0$.
Thus every retained hypothesis holds, with a strictly negative conclusion.
\end{proof}

\paragraph{The obstruction persists for positive infinite sequences.}
Take instead $g_n=1$ for all $n\ge0$. This is a strictly positive
log-concave sequence. For each of the four parameter values above,
the product formula gives
\[
 \lim_{n\to\infty}\left|
 \frac{1/\Gq(\nu+n+1)}{1/\Gq(\nu+n)}\right|=1-q,
\]
so the radius of convergence is $1/(1-q)=625/369$.
The constant coefficient of $\Delta_g$ is the same negative number
as in the theorem. Analyticity at $x=0$ therefore gives
$\Delta_g<0$ for all sufficiently small positive $x$.
The refutation thus also applies with no zero coefficients and with
$x$ strictly inside the convergence interval.

\paragraph{An additional finite check.}
For $g_0=g_1=1$ and $g_n=0$ for $n\ge2$, the same calculation gives
\[
 \frac{\Delta_g}{AB}
 =-\frac{95}{576}
  +\frac{3051749269}{3176712000}x
  +\frac{86127716}{67284525}x^2.
\]
At $x=1/100$ this equals
$-15036888769351/96889716000000<0$.
The accompanying exact-arithmetic verifier checks this polynomial
using the four shifted q-gamma recurrences.

\paragraph{Scope and status.}
The parameter restrictions above were checked in both the current
arXiv source and the full published paper, rather than inferred from
MathDB's abbreviated theorem references. The failure occurs inside
the retained negative-$\mu$ range and does not rely on a pole or on
setting $x=0$. It identifies a necessary repair to the domain of the
extension. A current literature and public-claim audit found no
matching full solution claim; its sources and limitations are recorded
in the accompanying audit. This manuscript has not received external
peer review.

\begin{thebibliography}{9}
\bibitem{KK}
S.~I. Kalmykov and D.~B. Karp,
\emph{Inequalities for series in q-shifted factorials and q-gamma
functions}, J. Math. Anal. Appl. \textbf{460} (2018), 332--351.
\href{https://doi.org/10.1016/j.jmaa.2017.11.041}
{doi:10.1016/j.jmaa.2017.11.041}.
Preprint: \href{https://arxiv.org/abs/1702.03788}{arXiv:1702.03788}.
\end{thebibliography}

\clearpage
\paragraph{Task 2: resolve the nonnegative-parameter version.}
Determine whether the full value and coefficient conclusions hold for
every $\alpha,\beta>0$ when $\mu\ge0$ and $(g_n)$ is doubly positive.
Begin by checking later claims for this restricted version.

\paragraph{Task 3: determine an admissible negative-parameter range.}
Find sharp additional restrictions on $\mu,\alpha,\beta$ below
$\mu=0$ that can hold uniformly over all doubly positive coefficient
sequences. Assess whether adding $\mu+\min(\alpha,\beta)\ge0$
is sufficient before attempting a general theorem.

\paragraph{Task 4: classify the two-term case.}
For $g_0,g_1>0$ and $g_n=0$ for $n\ge2$, classify when the
quadratic $\Delta_g(x)$ is nonnegative throughout $x>0$.
Compare the resulting domain with the conditions needed for
nonnegative coefficients.

\paragraph{Task 5: quantify persistence under positive tails.}
For the counterexample parameters, find the maximal interval adjacent
to zero on which $\Delta_g(x)<0$ for $g_n=1$.
Develop explicit bounds that remain valid for a specified class of
positive log-concave perturbations of this sequence.
\end{document}
